% L11b student companion deck: covered calls, protective collars (with TJUL), and
% the net Delta of a portfolio of shares and options.
%
% November 3, 2026: replaces the retired docs/ Notes. Same day, the cash secured put
% section was replaced by protective collars (instructor decision). The frame order mirrors the
% lecture notebook. The disclaimer is slide 2, by instructor preference. Notation
% follows the lecture: shares bought at S(0), strike K, call premium \mathcal{P}_{c};
% collar put at K_1 (premium \mathcal{P}_{1}) and call at K_2 (premium \mathcal{P}_{2}); profit at expiration \hat{P}(S(T)), accumulation factor
% \mathcal{D}_{T,0}(g_y) (always written in full), breakeven S^{\ddagger}, cost
% basis S^{\dagger}, contract multiplier m, number of contracts n, net Delta
% \Delta_{\text{net}}, options Delta \Delta_{\text{opts}}, composite Delta
% \Delta_{\mathcal{C}}. Maximum losses are listed as positive amounts.
\documentclass[aspectratio=169,11pt]{beamer}
\usepackage{vnslides}

\newcommand{\Cc}{\mathcal{C}}
\newcommand{\Dc}{\mathcal{D}}
\newcommand{\Pc}{\mathcal{P}}
\newcommand{\repo}{https://github.com/varnerlab/CHEME-5660-CourseRepository-Fall-2026/blob/main/lectures/week-11/L11b}

\title{\texorpdfstring{{\vnheavy L11b}}{L11b}: Covered Calls and Protective Collars}
\subtitle{CHEME 5660 Cornell University}
\author{Copyright Jeffrey Varner 2026. All Rights Reserved}

\begin{document}

\vntitlepage

\begin{frame}{Disclaimer and Risks}
  \vspace{0.6em}
  \textbf{\ink{Training only.}} This content is for training and informational
  purposes only. The teaching team does not make, give, or endorse any offer or
  solicitation to buy or sell securities or derivative products, or provide
  investment or trading advice or strategy.

  \vspace{1.1em}
  \textbf{\ink{Risk.}} Trading involves risk and is generally not recommended for
  those with limited resources, experience, or a low risk tolerance. Past
  performance does not guarantee future success.

  \vspace{1.1em}
  \textbf{\ink{Responsibility.}} You are responsible for your investment decisions
  based on your financial circumstances, objectives, risk tolerance, and
  liquidity needs.
\end{frame}

\begin{frame}{Objectives for Today}
  Today, we analyze two strategies that combine shares with options. What can
  each trade earn and lose, and how does its value change with the share price?

  \vspace{0.5em}
  {\large\textbf{\ink{Today's concepts}}}
  \begin{itemize}\setlength{\itemsep}{0.4em}
    \item \hilite{Covered calls}: shares plus a short call, two regimes,
      breakeven, maximum profit and loss.
    \item \hilite{Protective collars}: shares plus a long put and a short call,
      three regimes, breakeven, maximum profit and loss.
    \item \hilite{Net Delta}: the shares plus the weighted Deltas of the option
      legs, and the shares that make a position delta-neutral.
  \end{itemize}
\end{frame}

\begin{frame}{What You Need to Know}
  \vspace{0.3em}
  {\large\textbf{\ink{Core concepts}}}
  \begin{itemize}\setlength{\itemsep}{0.4em}
    \item Covered call profit per share:
      $\min(S(T),K)-\Dc_{T,0}(g_{y})\cdot(S(0)-\Pc_{c})$. The premium lowers the
      breakeven, and the call caps the profit at the strike.
    \item Collar profit per share, put at $K_{1}$, call at $K_{2}$:
      $S(T)+(K_{1}-S(T))^{+}-(S(T)-K_{2})^{+}-\Dc_{T,0}(g_{y})\cdot(S(0)+\Pc_{1}-\Pc_{2})$.
      The put sets a floor, and the call a cap.
    \item Net Delta: $\Delta_{\text{net}}=N_{S}+\Delta_{\text{opts}}$; holding
      $N_{S}=-\Delta_{\text{opts}}$ shares makes a position delta-neutral. A collar's
      put and call can only lower its net Delta from that of the shares.
  \end{itemize}

  \vspace{0.5em}
  {\large\textbf{\ink{For understanding}} \dimmed{(you won't reproduce these)}}
  \begin{itemize}\setlength{\itemsep}{0.3em}
    \item The derivation of the net Delta and the composite Delta table.
  \end{itemize}
\end{frame}

\begin{frame}{Lectures and Examples}
  \vspace{0.3em}
  {\normalsize\textbf{\ink{Lecture:}}\par\small
  \href{\repo/CHEME-5660-L11b-Lecture-CoveredCalls-Collars-Fall-2026.ipynb}{\dimmed{CHEME-5660-L11b-Lecture-CoveredCalls-Collars-Fall-2026.ipynb}}}

  \vspace{0.6em}
  {\normalsize\textbf{\ink{Examples:}}\par\small
  \href{\repo/CHEME-5660-L11b-Example-AnalysisCoveredCall-Fall-2026.ipynb}{\dimmed{CHEME-5660-L11b-Example-AnalysisCoveredCall-Fall-2026.ipynb}}\par
  \href{\repo/CHEME-5660-L11b-Example-StandardListed-SPY-ProtectiveCollar-Fall-2026.ipynb}{\dimmed{CHEME-5660-L11b-Example-StandardListed-SPY-ProtectiveCollar-Fall-2026.ipynb}}}

  \vspace{0.45em}
  {\normalsize\textbf{\ink{Review on your own:}}\par\small
  \href{\repo/CHEME-5660-L11b-Example-CoveredCallFund-SPY-Fall-2026.ipynb}{\dimmed{CHEME-5660-L11b-Example-CoveredCallFund-SPY-Fall-2026.ipynb}}}
\end{frame}

% ========== Company Profile: J.P. Morgan Asset Management
\begin{frame}{Company Profile: J.P. Morgan Asset Management}
  \href{https://am.jpmorgan.com/us/en/asset-management/adv/}{\hilite{J.P. Morgan Asset Management}}:
  the investment management business of JPMorganChase. It runs two large covered
  call funds, JEPI (S\&P 500) and JEPQ (Nasdaq-100):
  \begin{itemize}\setlength{\itemsep}{0.4em}
    \item \hilite{Shares plus sold options}: each fund owns U.S. large-cap stocks,
      sells options, and seeks monthly income from the premiums and dividends.
    \item \hilite{Equity-linked notes}: notes whose value is tied to stocks; as
      notes, they carry credit or counterparty risk.
    \item \hilite{Size}: about 45 (JEPI) and 44 (JEPQ) billion USD on September
      30, 2026.
  \end{itemize}
  \vspace{0.3em}
  \textbf{\ink{Connection:}} the premiums lower the breakeven, and the sold calls
  give up gains above the strike.
\end{frame}

% ========== Covered Calls
\begin{frame}{Covered Calls}
  A \hilite{covered call} owns shares and sells a call on them, collecting the
  premium. Investors use it when they expect the share price to stay flat or
  rise slightly.

  \vspace{0.6em}
  \textbf{\ink{Naked calls:}} selling a call without owning the shares carries
  unlimited risk if the share price rises. The shares in a covered call cover
  that obligation.

  \vspace{0.6em}
  Covered call ETFs, such as
  \href{https://www.globalxetfs.com/funds/qyld/}{\hilite{QYLD}} and
  \href{https://www.rexshares.com/fang/}{\hilite{FEPI}}, run this strategy for you.
  \slidenote{Video:}{\href{https://www.youtube.com/embed/Tid1rfVIc9c?rel=0}{\dimmed{SMB Capital, Simple \& Low Risk Investing Strategy}}}
\end{frame}

\begin{frame}{Covered Call: Profit at Expiration}
  Buy shares at $S(0)$ and sell one call per $m$ shares, with strike $K$, premium
  $\Pc_{c}$, and expiration $T$. The value at expiration minus today's cost grown
  to expiration is:
  \begin{align*}
    \hat{P}(S(T)) &= S(T) - \left(S(T) - K\right)^{+} - \Dc_{T,0}(g_{y})\cdot\left(S(0) - \Pc_{c}\right)\\
    &= \underbrace{\left(S(T) - \Dc_{T,0}(g_{y})\cdot S(0)\right)}_{\textstyle\text{shares}} + \underbrace{\Dc_{T,0}(g_{y})\cdot\Pc_{c} - \left(S(T) - K\right)^{+}}_{\textstyle\text{short call}}
  \end{align*}
  where $\Dc_{T,0}(g_{y})=e^{g_{y}T}$ is the accumulation factor and
  $\hat{P}(S(T))$ is in USD/share. For $n$ contracts, multiply by $mn$.
\end{frame}

\begin{frame}{Covered Call: Two Regimes}
  With no dividends and no early exercise, the call's payoff is zero when
  $S(T)\leq K$, so:
  {\small\begin{equation*}
    \boxed{
    \hat{P}(S(T)) = \begin{cases}
    \overbrace{K - \Dc_{T,0}(g_{y})\cdot S(0) + \Dc_{T,0}(g_{y})\cdot\Pc_{c}}^{\textstyle\text{realized gain}}, & \text{if } S(T) > K\;\text{(exercised)}\\
    \Dc_{T,0}(g_{y})\cdot\Pc_{c} + \underbrace{\left(S(T) - \Dc_{T,0}(g_{y})\cdot S(0)\right)}_{\textstyle\text{unrealized gain}}, & \text{if } S(T) \leq K\;\text{(expires worthless)}
    \end{cases}}
  \end{equation*}}
  \begin{itemize}\setlength{\itemsep}{0.3em}
    \item \hilite{Call exercised}: we sell the shares at $K$, so the whole profit
      is realized, and the call caps it.
    \item \hilite{Call expires worthless}: we keep the premium and the shares; the
      share gain or loss is unrealized.
  \end{itemize}
\end{frame}

\begin{frame}{Covered Call: Breakeven, Maximum Profit, and Maximum Loss}
  The maximum profit, breakeven $S^{\ddagger}$, and maximum loss are given below.
  Maximum losses are listed as positive amounts (USD/share):
  \begin{center}\small
    \begin{tabular}{ll}
      \toprule
      Maximum profit & $K-\Dc_{T,0}(g_{y})\cdot\left(S(0)-\Pc_{c}\right)$ for $S(T)\geq K$ \\
      Breakeven & $S^{\ddagger}=\Dc_{T,0}(g_{y})\cdot\left(S(0)-\Pc_{c}\right)$ \\
      Maximum loss & $\Dc_{T,0}(g_{y})\cdot\left(S(0)-\Pc_{c}\right)$ at $S(T)=0$ \\
      \bottomrule
    \end{tabular}
  \end{center}
  Compared with owning the shares alone, the premium lowers the breakeven, and
  the call gives up the gains above the strike. If we bought the shares earlier at
  a cost basis $S^{\dagger}$, the realized gain on the shares at exercise is
  $K-S^{\dagger}$.
  \slidenote{Example:}{\href{\repo/CHEME-5660-L11b-Example-AnalysisCoveredCall-Fall-2026.ipynb}{\dimmed{CHEME-5660-L11b-Example-AnalysisCoveredCall-Fall-2026.ipynb}}\hspace{0.8em}{\vnmedium Review on your own:}\hspace{0.65em}\href{\repo/CHEME-5660-L11b-Example-CoveredCallFund-SPY-Fall-2026.ipynb}{\dimmed{CHEME-5660-L11b-Example-CoveredCallFund-SPY-Fall-2026.ipynb}}}
\end{frame}

% ========== Protective Collars
\begin{frame}{Protective Collars}
  A \hilite{protective collar} adds a long put to a covered call: own the
  shares, buy a put at $K_{1}$ below the share price, and sell a call at
  $K_{2}>K_{1}$ above it, same expiration $T$. Per $m$ shares, with premiums
  $\Pc_{1}$ (put) and $\Pc_{2}$ (call), the profit at expiration (USD/share) is:
  \begin{equation*}
    \hat{P}(S(T)) = S(T) + \underbrace{\left(K_{1} - S(T)\right)^{+}}_{\textstyle\text{long put}} - \underbrace{\left(S(T) - K_{2}\right)^{+}}_{\textstyle\text{short call}} - \underbrace{\Dc_{T,0}(g_{y})\cdot\left(S(0) + \Pc_{1} - \Pc_{2}\right)}_{\textstyle\text{today's cost, grown}}
  \end{equation*}
  The put sets a floor under the shares, the call caps them, and the call
  premium pays for some or all of the put.
\end{frame}

\begin{frame}{Protective Collar: Three Regimes}
  With no dividends and no early exercise, at most one option pays at expiration:
  {\small\begin{equation*}
    \boxed{
    \hat{P}(S(T)) = \begin{cases}
    K_{1} - \Dc_{T,0}(g_{y})\cdot\left(S(0) + \Pc_{1} - \Pc_{2}\right), & \text{if } S(T) < K_{1}\;\text{(put exercised)}\\
    S(T) - \Dc_{T,0}(g_{y})\cdot\left(S(0) + \Pc_{1} - \Pc_{2}\right), & \text{if } K_{1} \leq S(T) \leq K_{2}\;\text{(neither)}\\
    K_{2} - \Dc_{T,0}(g_{y})\cdot\left(S(0) + \Pc_{1} - \Pc_{2}\right), & \text{if } S(T) > K_{2}\;\text{(call exercised)}
    \end{cases}}
  \end{equation*}}
  \begin{itemize}\setlength{\itemsep}{0.2em}
    \item \hilite{Put exercised}: we sell the shares at $K_{1}$; the loss is
      realized and stops growing.
    \item \hilite{Neither}: we keep the shares; the profit moves with the share price.
    \item \hilite{Call exercised}: we sell the shares at $K_{2}$; the profit is
      realized and capped.
  \end{itemize}
\end{frame}

\begin{frame}{Protective Collar: Breakeven, Maximum Profit, and Maximum Loss}
  The maximum profit, breakeven $S^{\ddagger}$, and maximum loss, a positive
  amount (USD/share), are:
  \begin{center}\small
    \begin{tabular}{ll}
      \toprule
      Maximum profit & $K_{2}-\Dc_{T,0}(g_{y})\cdot\left(S(0)+\Pc_{1}-\Pc_{2}\right)$ for $S(T)\geq K_{2}$ \\
      Breakeven & $S^{\ddagger}=\Dc_{T,0}(g_{y})\cdot\left(S(0)+\Pc_{1}-\Pc_{2}\right)$, if $K_{1}\leq S^{\ddagger}\leq K_{2}$ \\
      Maximum loss & $\Dc_{T,0}(g_{y})\cdot\left(S(0)+\Pc_{1}-\Pc_{2}\right)-K_{1}$ for $S(T)\leq K_{1}$${}^{a}$ \\
      \bottomrule
    \end{tabular}
  \end{center}
  {\footnotesize ${}^{a}$If $\Dc_{T,0}(g_{y})\cdot\left(S(0)+\Pc_{1}-\Pc_{2}\right)>K_{1}$; otherwise none.\par}
  \vspace{0.3em}
  When $\Pc_{1}=\Pc_{2}$, the call premium pays for the whole put: a
  \hilite{zero-cost collar}. Compared with the covered call, the put limits the
  loss, and it costs some of the call premium.
\end{frame}

\begin{frame}{TJUL: A Collar Inside a Fund}
  The \href{https://www.innovatoretfs.com/etf/?ticker=TJUL}{\hilite{Innovator Equity Defined Protection ETF (TJUL)}}
  holds a collar-like package of SPY options from July 1, 2025 to June 30, 2027.
  Its holdings on September 30, 2026, all expiring June 30, 2027:
  \begin{itemize}\setlength{\itemsep}{0.3em}
    \item \hilite{Purchased call, strike 11.09}: deep in the money, it moves almost
      one for one with SPY and stands in for the shares (without dividends).
    \item \hilite{Purchased put, strike 616.16}: close to the starting SPY level of
      617.85, the floor.
    \item \hilite{Sold call, strike 718.63}: sets the cap.
  \end{itemize}
  \vspace{0.3em}
  It seeks a 100\% buffer and a 13.61\% cap over the whole period, before its
  0.79\% expense ratio, for holders from first to last day. The cap is on the
  fund's value, not the call's distance above 617.85.
  \slidenote{Example:}{\href{\repo/CHEME-5660-L11b-Example-StandardListed-SPY-ProtectiveCollar-Fall-2026.ipynb}{\dimmed{CHEME-5660-L11b-Example-StandardListed-SPY-ProtectiveCollar-Fall-2026.ipynb}}\hspace{0.8em}{\vnmedium Source:}\hspace{0.65em}\href{https://www.innovatoretfs.com/pdf/tjul_factsheet.pdf}{\dimmed{TJUL fact sheet, September 30, 2026}}}
\end{frame}

% ========== Delta of Covered Calls and Collars
\begin{frame}{Delta of Covered Calls and Collars}
  Previously, we used Delta as an estimate of the probability of exercise. Its
  other use is \hilite{hedging}: offsetting potential losses with an opposite
  position in a related \hilite{counter asset}, often the underlying shares.

  \vspace{0.5em}
  Hold cash $N_{C}$ (USD), $N_{S}$ shares, and the legs of a composite contract
  $\Cc$. Leg $i$ has orientation $\theta_{i}$ ($1$ long, $-1$ short), $n_{i}$
  contracts, multiplier $m$, and premium $\Pc_{i}(S,t)$ (USD/share). The wealth
  at time $t$ is:
  \begin{equation*}
    W_{t}=N_{C}+N_{S}\,S+m\sum_{i\in\Cc}\theta_{i}n_{i}\Pc_{i}(S,t)
  \end{equation*}
\end{frame}

\begin{frame}{The Net Delta of a Portfolio}
  The \hilite{net Delta} is the sensitivity of the wealth to the share price:
  \begin{align*}
    \Delta_{\text{net}}(S,t)&=\frac{\partial W_{t}}{\partial S} = \underbrace{\frac{\partial N_C}{\partial S}}_{=0} + \underbrace{\frac{\partial (N_S S)}{\partial S}}_{=N_S} + m\sum_{i\in\Cc} \theta_i n_i \underbrace{\frac{\partial \Pc_i}{\partial S}}_{=\Delta_i(S,t)}\\
    &= N_S + \underbrace{m\sum_{i\in\Cc} \theta_i n_i \Delta_i(S,t)}_{\textstyle= \Delta_{\text{opts}}(S,t)}
  \end{align*}
  The options contribute $\Delta_{\text{opts}}=m\,\Delta_{\Cc}$, where
  $\Delta_{\Cc}(S,t)=\sum_{i\in\Cc}\theta_{i}n_{i}\Delta_{i}(S,t)$ is the
  \hilite{composite Delta}.
\end{frame}

\begin{frame}{Delta-Neutral Positions}
  A position is \hilite{delta-neutral} if $\Delta_{\text{net}}(S,t)=0$, which we
  get by holding $N_{S}=-\Delta_{\text{opts}}(S,t)$ shares. For example, short
  $10$ calls at $\Delta=0.6$ with $m=100$ give:
  \begin{equation*}
    \Delta_{\text{opts}}=100\times(-1)\times 10\times 0.6=-600
    \quad\Longrightarrow\quad N_{S}=600 \text{ shares}
  \end{equation*}
  If leg $i$ is exercised or assigned at time $\tau\le T$, set $\Delta_{i}=0$ for
  $t\geq\tau$ and update the shares:
  \begin{equation*}
    N_S \leftarrow N_S + \theta_i m n_i \;\;\text{(call)}\qquad
    N_S \leftarrow N_S - \theta_i m n_i \;\;\text{(put)}
  \end{equation*}
\end{frame}

\begin{frame}{Covered Call and Collar Delta}
  A covered call holds $N_{S}=mn$ shares and is short $n$ calls with Delta
  $\Delta_{c}$; a collar adds $n$ long puts with Delta $\Delta_{p}$ (put at
  $K_{1}$, call at $K_{2}$). The cash does not contribute:
  \begin{align*}
    \Delta_{\text{net}}(S,t) &= mn\left(1-\Delta_{c}(S,t)\right) && \text{covered call}\\
    \Delta_{\text{net}}(S,t) &= mn\left(1+\Delta_{p}(S,t)-\Delta_{c}(S,t)\right) && \text{collar}
  \end{align*}
  Since $0\leq\Delta_{c}\leq 1$ and $-1\leq\Delta_{p}\leq 0$, the put and the
  call can only lower the collar's net Delta from $mn$. Delta-neutral cases:
  \begin{itemize}\setlength{\itemsep}{0.2em}
    \item \hilite{Covered call}: $\Delta_{c}\approx 1$, the call deep in the money
      ($S(t)\gg K$).
    \item \hilite{Collar}: beyond either strike, $S(t)\ll K_{1}$ ($\Delta_{p}\approx-1$)
      or $S(t)\gg K_{2}$ ($\Delta_{c}\approx1$).
  \end{itemize}
\end{frame}

\begin{frame}{Delta of the Common Composite Contracts}
  Choose the strikes by Delta, a common rule of thumb: short legs at $\pm 0.25$,
  protecting long legs at $\pm 0.10$, the straddle at the money (about $\pm 0.5$). At
  entry, for one contract ($m=100$):\par
  {\centering\footnotesize\setlength{\tabcolsep}{5pt}
    \begin{tabular}{lll}
      \toprule
      Composite contract & Composite Delta, $\Delta_{\Cc}$ & Shares to hedge, $N_{S}$ \\
      \midrule
      Short put & $-(-0.25)=0.25$ & $-25$ \\
      Short call & $-0.25$ & $+25$ \\
      Put credit spread & $-0.10-(-0.25)=0.15$ & $-15$ \\
      Call credit spread & $-0.25+0.10=-0.15$ & $+15$ \\
      Short straddle & $-(-0.5)-0.5\approx0$ & $\approx0$ \\
      Short strangle & $-(-0.25)-0.25=0$ & $0$ \\
      Protected strangle & $-(-0.25)-0.25+0.10=0.10$ & $-10$ \\
      Iron condor & $-0.10-(-0.25)-0.25+0.10=0$ & $0$ \\
      \bottomrule
    \end{tabular}\par}
  \vspace{0.3em}
  {\small A Delta of zero is not zero risk: Gamma moves the Delta as the share
  price changes, so a delta-neutral position must update $N_{S}$.}
\end{frame}

\begin{frame}{Summary}
  \vspace{0.3em}
  We analyzed covered calls and protective collars, two strategies that combine
  shares with options, and computed the net Delta of shares plus options.

  \vspace{0.4em}
  {\large\textbf{\ink{Key takeaways}}}
  \begin{itemize}\setlength{\itemsep}{0.4em}
    \item \hilite{Covered calls}: the call caps the profit at the strike, and the
      premium lowers the breakeven.
    \item \hilite{Protective collars}: the put sets a floor and the call a cap;
      in between, the position moves with the shares.
    \item \hilite{Net Delta}: the shares plus the weighted Deltas of the option
      legs; a collar's put and call can only lower it from that of the shares.
  \end{itemize}
  \vspace{0.3em}
  These strategies trade some of the upside of the shares for income or for
  protection.
\end{frame}

\end{document}
